\PassOptionsToPackage{table}{xcolor}
\documentclass[pdflatex,iicol,sn-mathphys-ay]{sn-jnl}

\usepackage{graphicx}
\usepackage{amsmath}
\usepackage{amssymb}
\usepackage{comment}
\usepackage{algpseudocode}
\usepackage[linesnumbered,ruled,vlined]{algorithm2e}
\usepackage{multirow}
\usepackage[most]{tcolorbox}
\usepackage{float}
\usepackage{ulem}
\usepackage{xcolor}
\usepackage{subcaption}
\usepackage{stmaryrd}
\usepackage{tabularx}
\usepackage{booktabs}
\usepackage{makecell}
\usepackage{tikz}
\usepackage{array}
\usepackage{setspace}
\usepackage{pifont}

\newcommand{\cmark}{\ding{51}}
\newcommand{\xmark}{\ding{55}}
\newcommand{\omark}{\ding{108}}

\definecolor{darkblue}{rgb}{0.15,0.15,0.55}
\definecolor{lightgrey}{rgb}{0.75,0.75,0.75}
\newcommand{\qedm}{\hfill $\square$}
\DeclareUnicodeCharacter{25B3}{\ensuremath{\triangle}}
\newcommand\ny[1]{{\color{magenta}{#1}}}
\newcommand\sj[1]{{\color{violet}{#1}}}
\newcommand\sh[1]{{\color{orange}{#1}}}
\newcommand\gy[1]{{\color{purple}{#1}}}
\newcommand\cp[1]{{\color{blue}{#1}}}
\newcommand\jh[1]{{\color{green}{#1}}}

\hyphenation{op-tical net-works semi-conduc-tor}
\raggedbottom

\begin{document}

\title[Determining When and What to Ask in Robot Task Execution]{Determining When and What to Ask in Robot Task Execution}

\author*[1]{\fnm{Changmin} \sur{Park}}\email{cmp9877@gmail.com}
\author[1]{\fnm{Seongah} \sur{Park}}
\author[1]{\fnm{Nayoung} \sur{Kim}}
\author[1]{\fnm{Sungjin} \sur{Park}}
\author[1]{\fnm{Geonyeong} \sur{Ko}}
\author[2]{\fnm{Junho} \sur{Lee}}
\author[2]{\fnm{Changjoo} \sur{Nam}}

\affil*[1]{\orgname{Sogang University},
  \orgaddress{\city{Seoul}, \country{Republic of Korea}}}
\affil[2]{\orgname{Institution2},
  \orgaddress{\city{Seoul}, \country{Republic of Korea}}}

\abstract{We present a method for deciding when a robot should request external information and what it should ask during task execution under state uncertainty. Frequent queries incur interaction costs, while actions based on uncertain state facts can cause task failure. Our method combines online symbolic planning with a weighted belief over reachable successor states. Belief confidence determines whether to request state verification, and expected entropy reduction determines which Boolean state fact to query. The response updates the belief before the robot commits a symbolic state for subsequent planning. We evaluate the method in WasteSorting and TomatoHarvesting with oracle responses in simulation and VLM responses during physical robot execution. Confidence-based timing supports high task success across different query selection strategies, and entropy gain reduces query count at comparable success rates. In oracle simulation, our method achieves higher success with fewer queries than Query-Action. With VLM feedback, our method achieves higher success than KnowNo in both domains and requests more queries. These results show the distinct contributions of query timing and query selection to successful task execution.}

\keywords{Planning under uncertainty, online symbolic planning, robotic systems}

\maketitle

\section{Introduction}~\label{intro}

Robots operating in real-world environments often have incomplete or unreliable information about the current task state due to perception errors~\cite{kaipa2016addressing}, execution failures~\cite{kaelbling1998planning}, and unexpected environmental changes~\cite{kurniawati2016online}. If such uncertainty is not resolved, the robot may make incorrect decisions and fail the task. External information from an expert can help resolve this uncertainty~\cite{dragan2013legibility, goodrich2008human}, but frequent queries can incur substantial costs in time and workload~\cite{chen2014human}. The robot must therefore determine \textit{when} external information is necessary and \textit{what} information should be requested.

Recent uncertainty-aware systems selectively request external feedback during task execution. Existing approaches commonly trigger queries based on prediction confidence, including conformal prediction~\cite{ren2023robots, liang2024introspective, mullen2025lbap}, or on the expected effect of a query on task reward~\cite{burks2021collaborative, anand2025adaptive}. Prediction confidence measures uncertainty over candidate decisions, and expected task reward measures the value of a query for task execution. Neither criterion directly requires confidence in the current symbolic state to reach a specified level. Even if a single action is selected, the state facts required for execution can remain uncertain. We determine query timing from confidence in the symbolic state and query selection from expected entropy reduction over the remaining state hypotheses.

In this paper, we present a method that constructs a weighted belief over reachable successor states after each action and observation. We formulate decision making as a partially observable Markov decision process (POMDP)~\cite{kaelbling1998planning} and represent task knowledge using a STRIPS-style symbolic model~\cite{davis1993knowledge,fikes1971strips}. The belief is updated after each action and observation, and its uncertainty is used to determine when external information is needed. If a query is triggered, the robot selects the state fact expected to reduce uncertainty the most and asks an external expert to verify it. The response is incorporated into the belief, which is then used for subsequent planning and execution.


We evaluate the method in tomato harvesting and waste sorting, which exhibit different forms of state uncertainty. Using oracle feedback, we analyze the effects of query timing and query selection and compare against conformal-prediction-based and reward-based strategies. We then evaluate non-oracle feedback from a VLM and human participants and compare against KnowNo~\cite{ren2023robots}. The planned interviews examine how participants interpret robot questions and choose responses.


The main contributions of this work are as follows.

\begin{itemize}

\item We determine when to query based on uncertainty in the current task state and show that query timing has a substantial effect on task success.

\item We prioritize state information that is expected to reduce uncertainty and show that query selection affects the number of required queries.

\item We show higher task success with fewer queries than Query-Action in oracle simulation and higher success with additional queries than KnowNo during physical execution with VLM feedback.

\end{itemize}

\input{section/02_related}
\section{Preliminaries}~\label{preliminaries}
\subsection{Symbolic State and Action Representation}~\label{preliminaries:STRIPS}
The objective of planning is to find a transition sequence that transforms the initial state into a goal state. The planning problem is defined as a tuple $\mathcal{P} = \langle \mathcal{F}, \mathcal{A}, s_0, \mathcal{G} \rangle$, where $\mathcal{F}$ is a finite set of atoms, $\mathcal{A}$ is a set of actions, $s_0 \subseteq \mathcal{F}$ is the initial state, and $\mathcal{G} \subseteq \mathcal{F}$ is the goal condition. A state $s$ is represented as a subset of $\mathcal{F}$. Each action $a \in \mathcal{A}$ is defined as a tuple $a = \langle \mathrm{pre}(a), \mathrm{add}(a), \mathrm{del}(a) \rangle$, where $\mathrm{pre}(a), \mathrm{add}(a), \mathrm{del}(a) \subseteq \mathcal{F}$ denote the preconditions, add lists, and delete lists of the action, respectively. 
% We can say that $a$ is applicable in $s$ if and only if $\mathrm{pre}(a) \subseteq s$, denoted as $a \llbracket s \rrbracket$. 
An action $a$ is said to be applicable in state $s$ if and only if $\mathrm{pre}(a) \subseteq s$, denoted as $a \llbracket s \rrbracket$.
The state transition is defined as $s \xrightarrow{a} s^\prime$, where $s^\prime = (s \setminus \mathrm{del}(a)) \cup \mathrm{add}(a).$

A plan is a sequence of actions $\pi = (a_1, \dots, a_T)$ such that the goal condition is satisfied after sequentially applying the actions starting from the initial state, i.e.,
\[
s_0 \xrightarrow{a_1} s_1 \xrightarrow{a_2} \cdots \xrightarrow{a_T} s_T
\quad \text{and} \quad \mathcal{G} \subseteq s_T.
\]

\subsection{POMDP}~\label{preliminaries:POMDPs}
\begin{figure}[t]
\captionsetup{skip=0pt}
    \centering
    % \includegraphics[width=0.5\linewidth]
    \includegraphics[width=0.85\linewidth]
    {figures/fig_pomdp_preliminaries.pdf}
    \caption{Graphical illustration of a POMDP. At each time step $t$, the agent selects an action $a_t$ based on its current belief over the latent state $s_t$. The environment transitions to the next state $s_{t+1}$ and generates a reward $r_t$. Since the true state is not directly observable, the agent receives an observation $o_{t+1}$.}
    \label{fig:pomdp}
    % \vspace{-10pt}
\end{figure} 


A POMDP is defined as a tuple $\langle \mathcal{S}, \mathcal{A}, T, R, \Omega, O, \gamma \rangle$, where $\mathcal{S}$ is the state space, $\mathcal{A}$ is the action space, and $\Omega$ is the observation space. The transition model $T(s^\prime \mid s, a)$ specifies the probability of reaching the next state $s^\prime$ after taking action $a$ in current state $s$. The observation model $O(o \mid s^\prime, a)$ defines the probability of receiving observation $o$ after arriving at state $s^\prime$ by executing action $a$. The reward function $R(s,a)$ gives a scalar value to each transition, and $\gamma \in (0,1]$ is the discount factor.

Since the true state is not directly observable, the agent maintains a belief $b$, which is a probability distribution over $\mathcal{S}$ such that $b(s) = \Pr(s)$. After executing an action $a$ and receiving an observation $o$, the belief is updated using a Bayes filter.
\[
b^\prime(s^\prime) = \eta \, O(o \mid s^\prime,a) \sum_{s \in \mathcal{S}} T(s^\prime \mid s,a) \, b(s),
\]
where $\eta$ is a normalization constant.

A policy $\mu$ maps beliefs to actions. The objective of a POMDP is to find a policy that maximizes the expected cumulative discounted reward, defined as $\mathbb{E}\left[\sum_{t=0}^{\infty} \gamma^t r_t \right]$.


\subsection{POMCP}~\label{preliminaries:POMCP}
Partially Observable Monte Carlo Planning (POMCP) performs Monte Carlo tree search over action-observation histories $h = \langle a_1 o_1 \dots a_t o_t \rangle$, where each node in the search tree corresponds to a history. Action selection follows the UCB rule.
\begin{equation}~\label{eq:ucb}
a^* = \arg\max_a \left( V(ha) + c \sqrt{\frac{\log N(h)}{N(ha)}} \right),
\end{equation}
where $V(ha)$ is the estimated value of $a$ after $h$, $N(h)$ and $N(ha)$ denote visit counts, and $c$ is an exploration constant.

Each simulation samples the next state, observation, and reward from a generative model $(s^\prime, o, r) \sim \mathcal{G}(s,a)$. The return is computed recursively as $R = r + \gamma R^\prime$, where $R^\prime$ is the return from subsequent steps. The value estimate is updated using an incremental average:
\[
V(ha) \leftarrow V(ha) + \frac{R - V(ha)}{N(ha)}.
\]

The belief associated with a history $h$ is represented implicitly as a set of particles, $B(h) = \{ s \mid s \text{ is visited at } h \}$, which approximates the posterior over states without explicit belief updates.

\SetKwProg{Proc}{procedure}{}{}
\SetKwFunction{Search}{SEARCH}
\SetKwFunction{SampleParticle}{SAMPLE-PARTICLE}
\SetKwFunction{Simulate}{SIMULATE}
\SetKwFunction{Timeout}{TIMEOUT}
\SetKwFunction{Maxstep}{MAX-STEP}
\SetKwFunction{SearchBest}{SEARCH-BEST}
\SetKwFunction{UCB}{UCB}
\SetKwFunction{Rollout}{ROLLOUT}
\SetKwFunction{UpdateBelief}{UPDATE-BELIEF}
\SetKwFunction{ComputeConfidence}{CONF}
\SetKwFunction{SelectBestHypothesis}{SELECT-FACT}
\SetKwFunction{QuerySource}{QUERY}
\SetKwFunction{FilterBelief}{FILTER}

\section{Method}~\label{method:method}
To determine when and what to query the user, we maintain a particle-based belief state to quantify uncertainty and identify facts that distinguish possible task states. However, particle sampling in large state spaces is computationally expensive, and unconstrained sampling may produce states that violate symbolic constraints. To address these challenges, we propose a frontier-based particle sampling approach that uses symbolic action models to maintain a constraint-consistent belief for query selection.

\subsection{Planning and Querying}~\label{method:overview}\label{method:planning}\label{method:PM}
We maintain a knowledge state $s_t^{K}$ for symbolic planning and execution. For online action selection, we employ POMCP~\cite{silver2010monte}, restricting tree expansion and rollout to symbolically applicable actions. At each step, the planner selects an action $a_t\in\mathcal{A}$ applicable to $s_t^{K}$, denoted by $a_t\llbracket s_t^{K}\rrbracket$. Following execution and observation $o_{t+1}$, we construct a weighted belief $\mathcal{B}_{t+1}$ by sampling from the reachable-state frontier defined by $(s_t^{K},a_t)$. This restricts belief construction to symbolically valid successors while capturing uncertainty in the action outcome. The frontier approximates the next state distribution from the committed state $s_t^{K}$. It does not propagate every hypothesis from the preceding belief, so uncertainty discarded at commitment is not retained across steps.

\subsubsection{When to Query}~\label{method:when}\label{method:belief_rep}
% TODO reference qu2023stochastic is absent from ref.bib. Restore the citation after the source is identified.
To determine when to query, we assess uncertainty about the current state through the entropy of the weighted belief. If several successor states remain plausible after an action and no single state is dominant, the robot cannot confidently identify its current state. In this situation, belief entropy is high because the weight is distributed across several possible states. We use this property of entropy to compute confidence from normalized belief entropy, where lower confidence indicates greater state uncertainty.

We observe that a symbolic action usually changes only part of the state. Based on this observation, we define the \textit{frontier} of reachable successor states as follows.
\begin{equation}
\mathcal{F}_{t+1}
=
\mathrm{Frontier}(s_t^{K},a_t)
=
\{s_{t+1}^{1},\ldots,s_{t+1}^{j}\}.
\label{eq:frontier}
\end{equation}
For example, an action to pick a tomato can change which object the robot holds but preserves the robot location. Uncertain outcomes within the frontier are mutually exclusive, and no two successor states are identical. We restrict particle sampling to this frontier instead of the entire symbolic state space.

We construct the weighted belief by assigning a normalized weight $w_k$ to each of the $j$ sampled frontier states $s_{t+1}^{k}$.

\begin{equation}
\mathcal{B}_{t+1}
=
\{(s_{t+1}^{k},w_k)\}_{k=1}^{j},
\qquad
\sum_{k=1}^{j}w_k=1.
\label{eq:symbolic_belief}
\end{equation}
We compute these weights by normalizing the product of the transition probability and observation likelihood for each state. This update uses the point-mass prior $b_t(s)=\mathbf{1}[s=s_t^{K}]$ in the Bayes filter of Sec.~\ref{preliminaries:POMDPs}.
\begin{equation}
\begin{aligned}
\bar{w}_{k}
&=T(s_{t+1}^{k}\mid s_t^{K},a_t)\,
O(o_{t+1}\mid s_{t+1}^{k},a_t),\\
w_k&=\frac{\bar{w}_k}{\sum_{\ell=1}^{j}\bar{w}_{\ell}}.
\end{aligned}
\label{eq:belief_update}
\end{equation}
If all $\bar{w}_k$ are zero, we set $w_k=1/j$.

This explicit weighted frontier is used for confidence estimation and query selection. POMCP maintains sampled states in the search tree to estimate action values. Both components use the same symbolic transition and observation models. After verification, the committed state determines action applicability and the observation branch retained for the next search. The root particles remain those stored in that branch, with the committed state as a fallback if the branch contains no particles. The frontier weights are not copied into the search tree.

We compute the entropy of the belief as
\[
\mathcal{H}(\mathcal{B}_{t+1})
=
-\sum_{k=1}^{j}w_k\log_2 w_k.
\]
Because the maximum entropy depends on the number of represented worlds, we normalize by $\log_2 j$ and define
\begin{equation}
\mathrm{Conf}(\mathcal{B}_{t+1})
=
1-\frac{\mathcal{H}(\mathcal{B}_{t+1})}{\log_2 j}.
\label{eq:compute_confidence}
\end{equation}
For $j=1$, we define $\mathrm{Conf}(\mathcal{B}_{t+1})=1$. The score approaches zero if the represented states have similar weights and approaches one if most of the belief weight lies on one state.

Given a predefined threshold $\tau\in[0,1]$, the robot initiates state verification if
\begin{equation}
\mathrm{Conf}(\mathcal{B}_{t+1})<\tau.
\label{eq:when_query}
\end{equation}
A higher threshold $\tau$ requires a more concentrated belief and can increase the number of verification queries. Because filtering changes both entropy and the number of remaining states, confidence need not increase after every response. We recompute confidence over the retained states after each query.

\begin{figure}[t]
\captionsetup{skip=0pt}
\centering
\includegraphics[width=0.98\linewidth]{figures/fig_overview.pdf}
\caption{Online planning with belief-guided state verification. After executing an action and receiving an observation, the planner constructs a weighted belief over reachable symbolic states. Belief confidence determines whether to enter the query loop, and expected entropy reduction determines which symbolic fact to verify. The most probable remaining state becomes the symbolic reference for the next planning step.}
\label{fig:pm}
\end{figure}

\subsubsection{What to Query}~\label{method:what}
If belief confidence falls below $\tau$, the robot triggers a query and must decide which fact to verify. Since verification eliminates state hypotheses that are inconsistent with the answer, we select the fact with the largest expected reduction in belief entropy.

Specifically, we define the query candidate set $\mathcal{Q}(\mathcal{B})$ by extracting atomic facts whose truth values differ across the remaining frontier states in the belief. The candidates include all such facts without an additional restriction to action preconditions or goal facts.
\[
\mathcal{Q}(\mathcal{B})
=
\bigcup_{(s,w)\in\mathcal{B}}s
\setminus
\bigcap_{(s,w)\in\mathcal{B}}s.
\]
For each candidate fact $q\in\mathcal{Q}(\mathcal{B})$, we partition the belief according to whether $q$ holds in each state.
\[
\begin{aligned}
\mathcal{B}^{q+}&=\{(s^k,w_k)\mid q\in s^k\},\\
\mathcal{B}^{q-}&=\{(s^k,w_k)\mid q\notin s^k\}.
\end{aligned}
\]
Let $P(q)=\sum_{k:q\in s^k}w_k$ and $P(\neg q)=1-P(q)$. Entropy within each subset is computed after conditionally normalizing its weights. The expected posterior entropy of asking about $q$ is
\[
\mathbb{E}[\mathcal{H}\mid q]
=
P(q)\mathcal{H}(\mathcal{B}^{q+})
+
P(\neg q)\mathcal{H}(\mathcal{B}^{q-}).
\]
The selected query is
\begin{equation}
q^*
=
\arg\min_{q\in\mathcal{Q}(\mathcal{B})}
\left[
P(q)\mathcal{H}(\mathcal{B}^{q+})
+
P(\neg q)\mathcal{H}(\mathcal{B}^{q-})
\right].
\label{eq:select_hypo}
\end{equation}
Eq.~\ref{eq:select_hypo} selects the fact $q^*$ that maximizes expected entropy reduction. Under the binary response model, this reduction equals the Boolean entropy of $P(q)$, so the criterion favors facts that divide belief mass evenly. The symbolic model defines the candidate facts, and the selection criterion ranks candidates by information gain without a guarantee of maximum task reward. The robot then converts this fact into a binary query about the current state. For example, $q^*=\texttt{ripe}(\text{tomato}_1)$ produces the question \textit{`Is tomato 1 ripe?'}. Based on the response $y\in\{\texttt{true},\texttt{false}\}$, the robot updates the belief by removing states inconsistent with the reported truth value.
\begin{equation}
\mathcal{B}
\leftarrow
\begin{cases}
\{(s^k,w_k)\mid q^*\in s^k\}, & y=\texttt{true},\\
\{(s^k,w_k)\mid q^*\notin s^k\}, & y=\texttt{false}.
\end{cases}
\label{eq:filter}
\end{equation}
The retained weights are renormalized, confidence is recomputed, and another fact is selected if confidence remains below $\tau$. If no state matches a response, the implementation preserves the belief before that response. The update treats accepted responses as correct Boolean observations. An incorrect response can therefore remove the true state, and the query loop cannot restore a removed hypothesis. The loop ends if the threshold is met or no differentiating fact remains. The robot then commits the maximum a posteriori state
\[
s_{t+1}^{K}
=
\arg\max_{s^k\in\mathcal{B}_{t+1}}w_k
\]
for subsequent symbolic planning.

Alg.~\ref{alg:online_planning} explains the overall online planning and querying procedure. The robot first initializes the belief, step counter, and action sequence (line~\ref{alg1_planning:initialize}). At each step, POMCP selects an action (line~\ref{alg1_planning:search}), and execution stops if no action is available (line~\ref{alg1_planning:dead_end}). Otherwise, the robot executes the action  (lines~\ref{alg1_planning:execute} and~\ref{alg1_planning:record}), then constructs the weighted frontier belief and computes confidence (lines~\ref{alg1_planning:belief_update} and~\ref{alg1_planning:compute_conf}). If confidence falls below $\tau$, the robot selects a candidate fact (line~\ref{alg1_planning:select_hypo}). If no candidate remains, the query loop ends (line~\ref{alg1_planning:no_candidate}). Otherwise, the robot obtains an external response and filters the belief (lines~\ref{alg1_planning:query} and~\ref{alg1_planning:filter}), then recomputes confidence (line~\ref{alg1_planning:recompute_conf}). After the query loop, the most probable remaining state becomes the next symbolic planning state (line~\ref{alg1_planning:commit}). The planner retains the matching search branch for the next step (line~\ref{alg1_planning:prune}). Appx.~\ref{appx:algorithms} provides the detailed POMCP search, simulation, and rollout procedures.

\begin{algorithm}[t]
\caption{Belief-Guided State Verification and Planning}
\label{alg:online_planning}
\DontPrintSemicolon
\textbf{Input}\quad initial state $s_0^{K}$, confidence threshold $\tau$, step limit $T_{\max}$\;
\textbf{Output}\quad executed action sequence $\pi$\;
\BlankLine
$\mathcal{B}\gets\{(s_0^{K},1)\}\qquad t\gets0\qquad \pi\gets()$\label{alg1_planning:initialize}\;
\While{$t<T_{\max}$}{
    $a_t\gets$\Search{$\mathcal{B}$}\label{alg1_planning:search}\;
    \If{$a_t=\emptyset$}{
        \textbf{break}\label{alg1_planning:dead_end}\tcp*{dead end}
    }
    execute $a_t$ and receive $o_{t+1}$\label{alg1_planning:execute}\;
    $\pi\gets\pi\mathbin{\|}a_t$\label{alg1_planning:record}\;
    $\mathcal{B}\gets$\UpdateBelief{$s_t^{K},a_t,o_{t+1}$}\label{alg1_planning:belief_update}\;
    $c\gets\mathrm{Conf}(\mathcal{B})$\label{alg1_planning:compute_conf}\;
    \While{$c<\tau$}{
        $q^*\gets$\SelectBestHypothesis{$\mathcal{B}$}\label{alg1_planning:select_hypo}\;
        \If{$q^*=\emptyset$}{
            \textbf{break}\label{alg1_planning:no_candidate}\tcp*{no differentiating fact}
        }
        $y\gets$\QuerySource{$q^*$}\label{alg1_planning:query}\;
        $\mathcal{B}\gets$\FilterBelief{$\mathcal{B},q^*,y$}\label{alg1_planning:filter}\;
        $c\gets\mathrm{Conf}(\mathcal{B})$\label{alg1_planning:recompute_conf}\;
    }
    $s_{t+1}^{K}\gets\arg\max_{s^k\in\mathcal{B}}w_k$\label{alg1_planning:commit}\;
    retain the search branch for $(a_t,s_{t+1}^{K})$\label{alg1_planning:prune}\;
    $t\gets t+1$\label{alg1_planning:advance}\;
}
\KwRet $\pi$\label{alg1_planning:return}\;
\end{algorithm}

\subsection{System Implementation}~\label{method:modules}
\begin{figure*}[t!]
\captionsetup{skip=0pt}
\centering
\includegraphics[width=0.95\linewidth]{figures/fig_dataflow.pdf}
\caption{Framework for robot execution of the proposed planning and querying method.}
\label{fig:data_flow}
\end{figure*}

We implement a modular framework to execute the proposed planning and querying method on a robot, as shown in Fig.~\ref{fig:data_flow}. The Knowledge Base (KB) stores symbolic facts and action models, and Sensing and Perception (SP) provides observations and numeric state estimates from sensor data. The Planning Manager (PM) uses this information to select robot actions, construct the frontier belief, and decide when and what to query. The Feedback Manager (FM) converts the selected fact into a question and the response into evidence for the belief update. The User Interface (UI) presents the questions and collects responses during human interaction.

\label{method:KB}
The KB partitions information into domain knowledge $\mathcal{I}^{D}$, static structured information $\mathcal{I}^{S}$, and local information $\mathcal{I}^{L}$ obtained during execution.
\[
\mathcal{KB}=\mathcal{I}^{D}\sqcup\mathcal{I}^{S}\sqcup\mathcal{I}^{L}.
\]
Domain knowledge contains action schemas and task rules. Static information includes invariant facts and fixed map or robot parameters. Local information contains perceptual facts and numeric fluents that are updated during execution. SP maps raw sensor input to these task-level observations. For example, RGB-D sensing can update object pose and ripeness estimates, while the symbolic observation likelihood determines how those observations weight the frontier states.

\label{method:action_execution}
Each symbolic action invokes a robot skill. Actions such as $\texttt{navigate}(\text{robot},l_1,l_2)$ and $\texttt{pick}(\text{robot},\text{tomato}_1)$ specify task-level intentions, while an off-the-shelf motion planner such as MoveIt~\cite{chitta2012moveit} uses numeric fluents from the KB to generate low-level motion. This separation allows the feedback policy to operate on task-level facts without depending on a particular motion-planning implementation.

\label{method:FM}
The FM maps the selected predicate to a predefined question template and converts the response to a Boolean value. During human interaction, the tablet UI displays the camera view, highlights the referenced object with a bounding box, and presents simple response options. Participants answer questions about the displayed state fact, and the planner selects the next robot action.

\section{Experiments}~\label{experiment}
We evaluate our method in simulation and on physical robots. The evaluation covers WasteSorting and TomatoHarvesting. We first compare the complete query policies in simulation with correct oracle responses. We then evaluate physical robot execution with responses from a vision-language model (VLM) and human participants. Ablation experiments examine the contributions of query timing and query selection. A threshold sweep examines the balance between task success and query count.

The oracle experiments include $200$ trials per domain and condition. The VLM experiments include $40$ trials per domain and method. The planned human evaluation is designed for five participants with $10$ trials per domain and method for a total of $40$ trials.

\begin{figure}[t!]
    %\captionsetup{skip=0pt} % \hspace{0.3cm}
    \begin{center}
  
    \begin{subfigure}[b]{0.47\textwidth}
        \captionsetup{skip=0pt}
	    \includegraphics[width=\textwidth]{figures/WasteSortingEnv.jpg}
        \caption{WasteSorting domain.}
        \label{fig:dom1}
    \end{subfigure}
    \quad
    \begin{subfigure}[b]{0.47\textwidth}
        \captionsetup{skip=0pt}
	    \includegraphics[width=\textwidth]{figures/fig_exp_setup2.pdf}
	    \caption{TomatoHarvesting domain.}
        \label{fig:dom2}
    \end{subfigure}
    \caption{Physical robot environments. (a) A fixed UR5e robot sorts objects from four waste categories into designated bins in a tabletop workspace. (b) A mobile manipulator harvests ripe tomatoes and discards rotten tomatoes.}
  \label{fig:env_setting}
  
  \end{center}
\end{figure}


\subsection{Experimental Setup}~\label{exp:setup}
The oracle experiments use simulated task execution and exact external responses. The VLM and human experiments use physical robots in laboratory environments. Figure~\ref{fig:env_setting} shows the physical setups. \textbf{WasteSorting} uses a fixed UR5e robot arm. \textbf{TomatoHarvesting} uses a Summit mobile platform equipped with a UR5e robot arm. The mobile platform uses an RGB-D camera for perception and LiDAR for localization. WasteSorting uses SAM3 for object detection~\cite{carion2025sam3segmentconcepts}. TomatoHarvesting uses YOLOv12 for object detection~\cite{tian2025yolov12}.


\subsubsection{Task Domains}~\label{exp:setup_domains}
The WasteSorting task requires the robot to place every object in the correct bin. The UR5e robot is positioned between the objects and recycling bins in Fig.~\ref{fig:dom1}. All objects are within reach. The robot can detect objects and then pick and place them. Each object belongs to one of four waste categories corresponding to general waste or recyclable paper or cans or plastic. The categories are initially unknown. The main sources of uncertainty in this domain are occlusion and noisy sensor observations. In simulation we randomize the initial object arrangement for each trial. A trial succeeds if the robot places all objects in the correct bins within the step limit. Each correct placement yields a task reward of $+5$. Completion of the task yields an additional reward of $+10$. A penalty of $-10$ applies to the third and each subsequent detection action in a consecutive sequence.

The TomatoHarvesting task requires the robot to harvest fresh ripe tomatoes and discard rotten tomatoes. Four tomatoes are distributed across two stems in Fig.~\ref{fig:dom2}. The robot navigates between the docking station and the two stems. The \texttt{detect} action distinguishes ripe from unripe tomatoes. After a tomato is picked the \texttt{scan} action distinguishes fresh from rotten tomatoes. The robot uses \texttt{place} for fresh tomatoes and \texttt{discard} for rotten tomatoes. Unripe tomatoes remain on the stem. The main source of uncertainty in this domain is noisy sensor observations. In simulation we randomize tomato locations across the two stems for each trial. A trial succeeds if all fresh ripe tomatoes are harvested and all rotten tomatoes are discarded within the step limit. Correct placement of a fresh tomato or disposal of a rotten tomato yields a task reward of $+10$. Consecutive \texttt{detect} actions at the same stem and consecutive \texttt{scan} actions on the same tomato incur a penalty of $-5$ from the second action onward.


\subsubsection{Implementation Details}~\label{exp:details}
Our POMCP planner uses $100$ simulations per planning step with a maximum simulation depth of $20$. The discount factor is $\gamma=0.2$ and the UCB exploration constant is $c=1.0$. The default confidence threshold is $\tau=0.8$. We use Gemma 4 31B IT with NVFP4 quantization~\cite{gemmateam2026gemma4} as the external expert in the VLM experiments. Appendix~\ref{appx:settings} provides the discount cutoff, effective search depth, baseline configurations, and trial limits.

\subsubsection{Evaluation Metrics}~\label{exp:metrics}
Success rate is the fraction of all trials that achieve the task goal within the trial limit. Query count is the number of external requests per trial. Query probability is the fraction of physical steps that trigger at least one query. Plan length is the number of physical actions executed per trial. Search time is the cumulative directly recorded search duration per trial. Oracle plan time includes POMCP search for Ours and Query-Action and language-model planning for KnowNo and IntroPlan. Planning time in the physical result figures and tables is a residual estimate after recorded execution and interaction durations are subtracted from total time, as specified in Appendix~\ref{appx:metric_details}. Interaction time is the cumulative time spent on external requests and response processing per trial. We report interaction time only for the VLM and human experiments. We note that all metrics except success rate use successful trials in the main comparisons. Appendix~\ref{appx:all_trial_queries} additionally reports oracle query counts over all attempted trials.

% Numerical reference is the manuscript summary in experiments_final/04_Total_Summary.
% Figures use the manuscript reference and final physical experiment summaries.
\subsection{System Evaluation}~\label{exp:sys_eval}

\begin{table*}[t]
\centering
\caption{Query timing and query selection strategies. Columns indicate query timing criteria and rows indicate query selection criteria. KnowNo and IntroPlan and Query-Action are the baselines. Fully-Random and Q1--Q2 and W1--W3 are the ablation configurations. Ours serves as the reference method for both comparisons. A dash denotes a combination not evaluated in this study.}
\label{tab:query_strategy_matrix}
\small
\newcommand{\queryhighlight}[2]{%
    \tikz[baseline=(marker.base)]{%
        \node[inner xsep=3pt,inner ysep=1pt] (marker) {\phantom{#2}};
        \fill[#1,opacity=0.7]
            ([xshift=-1pt,yshift=-0.5pt]marker.south west) --
            ([xshift=1pt,yshift=0.5pt]marker.south east) --
            ([xshift=2pt,yshift=-0.5pt]marker.north east) --
            ([xshift=0.5pt,yshift=0.5pt]marker.north west) -- cycle;
        \node[inner xsep=3pt,inner ysep=1pt] at (marker.center) {#2};
    }%
}
\begin{tabular}{lcccc}
\toprule
\shortstack{When $\rightarrow$\\What $\downarrow$} & Random & Conformal Prediction & Reward-based & Confidence \\
\midrule
Random & \queryhighlight{blue!25}{Fully-Random} & -- & -- & \queryhighlight{blue!25}{Q1} \\
Conformal Set & -- & \queryhighlight{yellow!65}{KnowNo, IntroPlan} & -- & -- \\
Reward-based & -- & -- & \queryhighlight{yellow!65}{Query-Action} & \queryhighlight{blue!25}{Q2} \\
Entropy Gain & \queryhighlight{blue!25}{W1} & \queryhighlight{blue!25}{W2} & \queryhighlight{blue!25}{W3} & \queryhighlight{green!35}{\textbf{Ours}} \\
\bottomrule
\end{tabular}
\end{table*}

Table~\ref{tab:query_strategy_matrix} summarizes the baseline and ablation configurations. The baseline comparison in Sec.~\ref{exp:sys_baseline} evaluates KnowNo, IntroPlan, and Query-Action against Ours. The baseline evaluation covers oracle responses in simulation and VLM and human responses during physical execution. The ablation study in Sec.~\ref{exp:policy_ablation} evaluates Fully-Random, Q1--Q2, and W1--W3 against Ours where W1--W3 examine query timing (when) and Q1--Q2 examine query selection (what). Finally the threshold study in Sec.~\ref{exp:sys_threshold} examines how the confidence threshold affects task success and query count. Both the ablation and threshold studies use oracle responses in simulation.



\subsubsection{Baseline Comparison}~\label{exp:sys_baseline}

\noindent\textbf{Oracle}. The oracle comparison includes three baselines with the query timing and selection strategies summarized in Table~\ref{tab:query_strategy_matrix}. KnowNo~\cite{ren2023robots} requests an action choice if its conformal prediction set does not specify a single ordinary action. IntroPlan~\cite{liang2024introspective} extends KnowNo and incorporates retrieved examples and explanations into candidate evaluation before it requests an action choice. Query-Action follows the standard POMDP approach to information gathering~\cite{kaelbling1998planning}. It includes Boolean state queries in the POMCP action space and selects them according to expected cumulative reward. Table~\ref{tab:baseline_results} and Fig.~\ref{fig:baseline_results} compare the complete planning and querying systems. The baselines differ in action selection, query criteria, and recorded planner settings. Appendix~\ref{appx:baseline_configs} specifies these settings, and Appendix~\ref{appx:oracle_budget_audit} examines the trial-budget difference. 

\begin{figure*}[!tp]
\centering
\begin{subfigure}[t]{0.247\textwidth}
\centering
\includegraphics[width=\linewidth,height=.16\textheight,keepaspectratio]{figures/exp_oracle_waste_success.pdf}
\caption{WasteSorting\\Success rate (\%).}
\label{fig:baseline_oracle_waste_success}
\end{subfigure}%
\hspace{0.004\textwidth}%
\begin{subfigure}[t]{0.247\textwidth}
\centering
\includegraphics[width=\linewidth,height=.16\textheight,keepaspectratio]{figures/exp_oracle_waste_queries.pdf}
\caption{WasteSorting\\Query count.}
\label{fig:baseline_oracle_waste_queries}
\end{subfigure}%
\hspace{0.004\textwidth}%
\begin{subfigure}[t]{0.247\textwidth}
\centering
\includegraphics[width=\linewidth,height=.16\textheight,keepaspectratio]{figures/exp_oracle_waste_query_probability.pdf}
\caption{WasteSorting\\Query probability (\%).}
\label{fig:baseline_oracle_waste_query_probability}
\end{subfigure}%
\hspace{0.004\textwidth}%
\begin{subfigure}[t]{0.247\textwidth}
\centering
\includegraphics[width=\linewidth,height=.16\textheight,keepaspectratio]{figures/exp_oracle_waste_plan_time.pdf}
\caption{WasteSorting\\Plan time (s).}
\label{fig:baseline_oracle_waste_plan_time}
\end{subfigure}
\par\smallskip
\begin{subfigure}[t]{0.247\textwidth}
\centering
\includegraphics[width=\linewidth,height=.16\textheight,keepaspectratio]{figures/exp_oracle_tomato_success.pdf}
\caption{TomatoHarvesting\\Success rate (\%).}
\label{fig:baseline_oracle_tomato_success}
\end{subfigure}%
\hspace{0.004\textwidth}%
\begin{subfigure}[t]{0.247\textwidth}
\centering
\includegraphics[width=\linewidth,height=.16\textheight,keepaspectratio]{figures/exp_oracle_tomato_queries.pdf}
\caption{TomatoHarvesting\\Query count.}
\label{fig:baseline_oracle_tomato_queries}
\end{subfigure}%
\hspace{0.004\textwidth}%
\begin{subfigure}[t]{0.247\textwidth}
\centering
\includegraphics[width=\linewidth,height=.16\textheight,keepaspectratio]{figures/exp_oracle_tomato_query_probability.pdf}
\caption{TomatoHarvesting\\Query probability (\%).}
\label{fig:baseline_oracle_tomato_query_probability}
\end{subfigure}%
\hspace{0.004\textwidth}%
\begin{subfigure}[t]{0.247\textwidth}
\centering
\includegraphics[width=\linewidth,height=.16\textheight,keepaspectratio]{figures/exp_oracle_tomato_plan_time.pdf}
\caption{TomatoHarvesting\\Plan time (s).}
\label{fig:baseline_oracle_tomato_plan_time}
\end{subfigure}
\caption{Oracle baseline comparison of success rate and query count and query probability and plan time. Our method achieves higher success with fewer queries and less plan time than Query-Action in both domains.}
\label{fig:baseline_results}
\end{figure*}


\begin{table*}[t]
\centering
\caption{Oracle baseline results by domain.}
\label{tab:baseline_results}
\small
\setlength{\tabcolsep}{4pt}
\begin{tabular}{llcccc}
\toprule
Domain & Method & Success (\%) & Query count & \shortstack{Query prob.(\%)} & Plan length \\
\midrule
WasteSorting & KnowNo & 48.50 & $2.48 \pm 1.42$ & $22.82 \pm 12.61$ & $10.78 \pm 1.35$ \\
 & IntroPlan & 51.50 & $6.83 \pm 1.53$ & $61.97 \pm 9.35$ & $10.99 \pm 1.49$ \\
 & Query-Action & 87.00 & $24.52 \pm 5.39$ & $87.05 \pm 5.28$ & $11.89 \pm 1.53$ \\
 & Ours & 100.00 & $7.83 \pm 2.16$ & $30.55 \pm 10.41$ & $10.35 \pm 0.88$ \\
\midrule
TomatoHarvesting & KnowNo & 43.00 & $2.45 \pm 1.06$ & $17.27 \pm 6.67$ & $14.01 \pm 1.07$ \\
 & IntroPlan & 40.00 & $8.96 \pm 2.04$ & $61.13 \pm 10.77$ & $14.61 \pm 1.46$ \\
 & Query-Action & 57.50 & $11.28 \pm 3.78$ & $51.68 \pm 10.45$ & $14.94 \pm 1.70$ \\
 & Ours & 98.00 & $9.16 \pm 1.90$ & $40.08 \pm 4.56$ & $13.97 \pm 1.15$ \\
\bottomrule
\end{tabular}
\end{table*}

As shown in Fig.~\ref{fig:baseline_results}, our method achieves the highest success rate in both domains. It completes all $200$ WasteSorting trials and $196$ of $200$ TomatoHarvesting trials. Our method achieves success rates of $100.00\%$ in WasteSorting and $98.00\%$ in TomatoHarvesting compared with $87.00\%$ and $57.50\%$ for Query-Action. The mean query counts are also lower at $7.83$ and $9.16$ compared with $24.52$ and $11.28$ for Query-Action. Two differences in planning may contribute to the lower success rate and higher query count of Query-Action. First, query actions expand the action space so POMCP must evaluate more alternatives within the same simulation budget. Second, expected task reward does not explicitly measure confidence in the current state. Our method selects queries at the observation stage and restricts POMCP planning to physical actions. These results show that a higher query count alone does not guarantee a higher success rate. Query-Action also counts queries toward the trial limit. Seven WasteSorting trials and $13$ TomatoHarvesting trials exhaust the $50$-decision budget, so the baseline comparison includes an effect of different execution limits.

Compared to KnowNo and IntroPlan, our method uses more queries. Their mean query counts are $2.48$ and $6.83$ in WasteSorting compared with $7.83$ for our method. In TomatoHarvesting, the corresponding counts are $2.45$ and $8.96$ compared with $9.16$ for our method. This difference in query count arises from how each method handles a query trigger. If a query is triggered, our method repeats queries until belief confidence reaches $\tau$. KnowNo and IntroPlan ask only one question per trigger. Therefore our method can use more queries despite a lower query probability than IntroPlan. However KnowNo and IntroPlan achieve success rates of only $48.50\%$ and $51.50\%$ in WasteSorting and $43.00\%$ and $40.00\%$ in TomatoHarvesting. Therefore a system is not necessarily better simply because it uses fewer queries.

Plan time is also lowest for Ours in both domains despite the additional queries compared with KnowNo and IntroPlan. In WasteSorting, mean plan time is $0.91$ seconds for Ours compared with $22.54$ seconds for KnowNo, $44.43$ seconds for IntroPlan, and $10.47$ seconds for Query-Action. In TomatoHarvesting, the corresponding values are $1.10$, $20.72$, $73.80$, and $5.55$ seconds. KnowNo and IntroPlan require language-model inference for candidate generation and scoring at each planning step. IntroPlan performs additional reasoning and candidate evaluation, which increases inference time compared with KnowNo. Ours evaluates query timing and selection from the belief state and searches only over physical actions. The smaller search space supports faster planning than Query-Action. In the evaluated implementations, search over symbolic actions reduces planning latency without repeated language-model action generation. % Therefore a higher query count does not necessarily imply a higher planning cost.

Our method combines confidence-based query timing with entropy gain query selection and achieves the highest success rate in both domains. The baseline comparisons suggest that effective query timing and selection are more important than query count alone.

\noindent\textbf{VLM}\label{exp:vlm_eval}. We evaluate whether the trends in task success and query behavior from the oracle experiments also hold in a real-world setting. We first prompt a VLM to act as an external expert and evaluate query policies on physical robots.\footnote{Appendix~\ref{appx:vlm_prompts} provides the VLM prompts.} We select KnowNo as the baseline because it achieves a high success rate relative to its query count in the oracle experiments. Each method has $40$ trials per domain. Figure~\ref{fig:vlm_results} presents the VLM results.

\begin{figure*}[!tp]
\centering
\begin{subfigure}[t]{0.247\textwidth}
\centering
\includegraphics[width=\linewidth,height=.20\textheight,keepaspectratio]{figures/exp_vlm_success.pdf}
\caption{Success rate (\%).}
\label{fig:vlm_vlm_success}
\end{subfigure}%
\hspace{0.004\textwidth}%
\begin{subfigure}[t]{0.247\textwidth}
\centering
\includegraphics[width=\linewidth,height=.20\textheight,keepaspectratio]{figures/exp_vlm_queries.pdf}
\caption{Query count.}
\label{fig:vlm_vlm_queries}
\end{subfigure}%
\hspace{0.004\textwidth}%
\begin{subfigure}[t]{0.247\textwidth}
\centering
\includegraphics[width=\linewidth,height=.20\textheight,keepaspectratio]{figures/exp_vlm_query_probability.pdf}
\caption{Query probability (\%).}
\label{fig:vlm_vlm_query_probability}
\end{subfigure}%
\hspace{0.004\textwidth}%
\begin{subfigure}[t]{0.247\textwidth}
\centering
\includegraphics[width=\linewidth,height=.20\textheight,keepaspectratio]{figures/exp_vlm_plan_length.pdf}
\caption{Plan length (actions).}
\label{fig:vlm_vlm_plan_length}
\end{subfigure}
\par\smallskip
\begin{subfigure}[t]{0.247\textwidth}
\centering
\includegraphics[width=\linewidth,height=.20\textheight,keepaspectratio]{figures/exp_vlm_total_time.pdf}
\caption{Total time (s).}
\label{fig:vlm_vlm_total_time}
\end{subfigure}%
\hspace{0.004\textwidth}%
\begin{subfigure}[t]{0.247\textwidth}
\centering
\includegraphics[width=\linewidth,height=.20\textheight,keepaspectratio]{figures/exp_vlm_interaction_time.pdf}
\caption{Interaction time (s).}
\label{fig:vlm_vlm_interaction_time}
\end{subfigure}%
\hspace{0.004\textwidth}%
\begin{subfigure}[t]{0.247\textwidth}
\centering
\includegraphics[width=\linewidth,height=.20\textheight,keepaspectratio]{figures/exp_vlm_planning_time.pdf}
\caption{Planning time (s).}
\label{fig:vlm_vlm_planning_time}
\end{subfigure}
\caption{Physical robot evaluation with VLM responses. Ours achieves $92.50\%$ success in both domains. Query probability is the fraction of physical steps with at least one query.}
\label{fig:vlm_results}
\end{figure*}


\begin{table*}[t]
\centering
\caption{Physical robot results with VLM and human responses. Parentheses give successful trials over all trials. Time is measured in seconds.}
\label{tab:physical_robot_results}
\small
\begin{tabular}{llccccc}
\toprule
Source & Method & Success (\%) & Queries & Total time & Interaction & Planning \\
\midrule
\multicolumn{7}{l}{WasteSorting} \\
VLM & KnowNo & 50.00 (20/40) & $3.15 \pm 0.88$ & $186.51 \pm 10.90$ & $15.46 \pm 4.45$ & $20.08 \pm 4.76$ \\
 & Ours & 92.50 (37/40) & $6.70 \pm 1.76$ & $193.69 \pm 13.52$ & $26.83 \pm 7.76$ & $1.11 \pm 0.05$ \\
Human & KnowNo & 16.67 (1/6) & $4.00 \pm \text{--}$ & $183.38 \pm \text{--}$ & $10.78 \pm \text{--}$ & $20.40 \pm \text{--}$ \\
 & Ours & 100.00 (6/6) & $6.00 \pm 1.10$ & $167.45 \pm 7.83$ & $15.65 \pm 4.90$ & $1.09 \pm 0.03$ \\
\midrule
\multicolumn{7}{l}{TomatoHarvesting} \\
VLM & KnowNo & 25.00 (10/40) & $7.20 \pm 1.87$ & $217.82 \pm 17.78$ & $44.99 \pm 13.10$ & $27.68 \pm 3.20$ \\
 & Ours & 92.50 (37/40) & $8.51 \pm 1.52$ & $180.39 \pm 11.18$ & $34.12 \pm 7.62$ & $1.48 \pm 0.19$ \\
Human & KnowNo & 33.33 (2/6) & $4.00 \pm 1.41$ & $181.89 \pm 3.38$ & $16.75 \pm 11.18$ & $25.40 \pm 3.56$ \\
 & Ours & 50.00 (3/6) & $9.67 \pm 1.53$ & $182.36 \pm 3.34$ & $27.00 \pm 2.64$ & $1.46 \pm 0.13$ \\
\bottomrule
\end{tabular}
\end{table*}

Ours achieves a success rate of $92.50\%$ in both domains with VLM responses. KnowNo achieves $50.00\%$ in WasteSorting and $25.00\%$ in TomatoHarvesting. Mean query counts are $6.70$ for Ours and $3.15$ for KnowNo in WasteSorting and $8.51$ and $7.20$ in TomatoHarvesting. Consistent with the oracle experiments, Ours requires less search time than KnowNo. Mean search time is $0.99$ seconds for Ours and $19.99$ seconds for KnowNo in WasteSorting and $1.38$ and $26.90$ seconds in TomatoHarvesting. In TomatoHarvesting, interaction time and total time are $44.99$ and $217.82$ seconds for KnowNo compared with $34.12$ and $180.39$ seconds for Ours. Notably, KnowNo uses only $15.43\%$ fewer queries than Ours in TomatoHarvesting with VLM responses compared with $73.22\%$ fewer queries in the oracle experiments. However the similar query counts do not translate into similar execution costs. KnowNo executes longer plans, with $16.10$ physical actions compared with $13.27$ for Ours. Interaction time and total time for KnowNo are $31.86\%$ and $20.75\%$ longer than for Ours respectively.

To understand the differences from the oracle results, we analyze failed trials and distinguish VLM response errors from baseline planner failures. In WasteSorting, all three failures of Ours involve VLM feedback. For example, the VLM rejects a can classification from perception and the robot places the object in the plastic bin. KnowNo has three VLM response failures ($15\%$) and $17$ planner failures ($85\%$). The planner failures occur because KnowNo selects actions without requesting assistance despite unmet action preconditions. In TomatoHarvesting, all three failures of Ours also involve VLM feedback. KnowNo has $18$ VLM failures ($60\%$) and $12$ planner failures ($40\%$). The planner failures occur after action selections without external assistance. In six trials, the VLM repeatedly selects detection despite feasible alternatives and the robot reaches the step limit. The remaining $12$ VLM failures involve incorrect action selections.

The WasteSorting failures show that a single predicted action can still require additional state verification before execution. In TomatoHarvesting, action selection requires the VLM to reason about robot location and object states to determine the next action. The spatial reasoning requirement may contribute to less effective action choices. Even successful trials can then require additional queries and longer plans and more total time. Therefore fewer queries do not necessarily reduce interaction time or total time. Ours asks the VLM to verify current state facts and assigns action selection to the planner. The VLM assesses current observations with less need to reason about future actions and outcomes. Performance may also depend on what the VLM is asked to judge.

\noindent\textbf{Human}\label{exp:human_eval}. We further examine whether the trends in task success and query behavior persist with human feedback. We provide five participants with domain information including the task goals and available actions. \textit{The provisional participant profile assumes two men and three women with a mean age of approximately $25$ years and an age range of $20$--$30$ years. These demographic values are placeholders and are not yet verified.} Participants use the tablet interface in Fig.~\ref{fig:human_tablet_interface} to answer robot questions. Participants choose responses based on their own judgment. We use the same robot platforms and task domains as in the VLM experiments. Each method has $10$ trials per domain. Figure~\ref{fig:human_results} presents the human evaluation results. Table~\ref{tab:physical_robot_results} reports the results of both the VLM and human evaluations.
% Replace the provisional demographics after participant data collection.


% Update the numerical results after the five-participant evaluation. Current values use six trials per domain and method.
% Planning times use the successful-trial values provided in the human experiment table.
Ours achieves a success rate of $100.00\%$ in WasteSorting and $50.00\%$ in TomatoHarvesting with human responses. KnowNo achieves $16.67\%$ and $33.33\%$ respectively. Mean query counts are $6.00$ for Ours and $4.00$ for KnowNo in WasteSorting and $9.67$ and $4.00$ in TomatoHarvesting. Mean planning time is $1.09$ seconds for Ours and $20.40$ seconds for KnowNo in WasteSorting and $1.46$ and $25.40$ seconds in TomatoHarvesting. In WasteSorting, interaction time and total time are $15.65$ and $167.45$ seconds for Ours compared with $10.78$ and $183.38$ seconds for KnowNo. In TomatoHarvesting, the corresponding times are $27.00$ and $182.36$ seconds for Ours compared with $16.75$ and $181.89$ seconds for KnowNo. KnowNo uses $33.33\%$ fewer queries than Ours in WasteSorting and $58.62\%$ fewer in TomatoHarvesting. However the lower query counts do not consistently reduce total time. Plan lengths are similar across methods, with $11.00$ physical actions for both methods in WasteSorting and $13.33$ for Ours and $13.00$ for KnowNo in TomatoHarvesting. Despite the additional queries, Ours requires $8.69\%$ less total time in WasteSorting and a similar total time in TomatoHarvesting.

\noindent\textbf{Comparison Across Feedback Sources}\label{exp:response_comparison}. Figure~\ref{fig:response_comparison} compares task success and query count across oracle, VLM, and human feedback. Ours achieves higher success rates and uses more queries than KnowNo in both domains across all three feedback sources. However the magnitude of the success advantage depends on the domain and feedback source. In WasteSorting, Ours maintains success rates of $100.00\%$, $92.50\%$, and $100.00\%$ with oracle, VLM, and human feedback respectively. In TomatoHarvesting, success decreases from $98.00\%$ with oracle feedback to $92.50\%$ with VLM feedback and $50.00\%$ with human feedback. The success advantage over KnowNo also narrows from $55.00$ percentage points with oracle feedback to $16.67$ percentage points with human feedback in TomatoHarvesting. Thus the relative advantage of Ours persists, but high success does not persist equally across domains and feedback sources.

Query counts also show differences across feedback sources. In TomatoHarvesting, Ours uses $9.16$, $8.51$, and $9.67$ queries with oracle, VLM, and human feedback respectively. KnowNo uses $2.45$, $7.20$, and $4.00$ queries. The query count of KnowNo increases substantially with VLM feedback, while the query count of Ours remains similar across feedback sources despite the lower success with human feedback. Therefore similar query counts do not imply similar task outcomes, and the number of queries alone does not explain the differences across feedback sources. These differences motivate the subsequent human interview analysis. We examine how participants interpret robot questions and what reasoning participants need to choose a response, with particular attention to question clarity and response effort.


\begin{figure}[t]
\centering
\includegraphics[width=\linewidth]{figures/human_tablet_interface.jpg}
\caption{Tablet interface used in the human evaluation. Participants view the robot camera feeds and select responses to robot queries.}
\label{fig:human_tablet_interface}
\end{figure}

\begin{figure*}[!tp]
\centering
\begin{subfigure}[t]{0.247\textwidth}
\centering
\includegraphics[width=\linewidth,height=.20\textheight,keepaspectratio]{figures/exp_human_success.pdf}
\caption{Success rate (\%).}
\label{fig:human_human_success}
\end{subfigure}%
\hspace{0.004\textwidth}%
\begin{subfigure}[t]{0.247\textwidth}
\centering
\includegraphics[width=\linewidth,height=.20\textheight,keepaspectratio]{figures/exp_human_queries.pdf}
\caption{Query count.}
\label{fig:human_human_queries}
\end{subfigure}%
\hspace{0.004\textwidth}%
\begin{subfigure}[t]{0.247\textwidth}
\centering
\includegraphics[width=\linewidth,height=.20\textheight,keepaspectratio]{figures/exp_human_query_probability.pdf}
\caption{Query probability (\%).}
\label{fig:human_human_query_probability}
\end{subfigure}%
\hspace{0.004\textwidth}%
\begin{subfigure}[t]{0.247\textwidth}
\centering
\includegraphics[width=\linewidth,height=.20\textheight,keepaspectratio]{figures/exp_human_plan_length.pdf}
\caption{Plan length (actions).}
\label{fig:human_human_plan_length}
\end{subfigure}
\par\smallskip
\begin{subfigure}[t]{0.247\textwidth}
\centering
\includegraphics[width=\linewidth,height=.20\textheight,keepaspectratio]{figures/exp_human_total_time.pdf}
\caption{Total time (s).}
\label{fig:human_human_total_time}
\end{subfigure}%
\hspace{0.004\textwidth}%
\begin{subfigure}[t]{0.247\textwidth}
\centering
\includegraphics[width=\linewidth,height=.20\textheight,keepaspectratio]{figures/exp_human_interaction_time.pdf}
\caption{Interaction time (s).}
\label{fig:human_human_interaction_time}
\end{subfigure}%
\hspace{0.004\textwidth}%
\begin{subfigure}[t]{0.247\textwidth}
\centering
\includegraphics[width=\linewidth,height=.20\textheight,keepaspectratio]{figures/exp_human_planning_time.pdf}
\caption{Planning time (s).}
\label{fig:human_human_planning_time}
\end{subfigure}
\caption{Physical robot evaluation with human responses. Ours completes six WasteSorting trials and three TomatoHarvesting trials. An asterisk denotes a single successful trial with no sample standard deviation.}
\label{fig:human_results}
\end{figure*}

\begin{figure*}[!tp]
\centering
\begin{subfigure}[t]{0.495\textwidth}
\centering
\includegraphics[width=\linewidth,height=.20\textheight,keepaspectratio]{figures/exp_response_waste_success.pdf}
\caption{WasteSorting. Success rate (\%).}
\label{fig:response_response_waste_success}
\end{subfigure}
\hfill
\begin{subfigure}[t]{0.495\textwidth}
\centering
\includegraphics[width=\linewidth,height=.20\textheight,keepaspectratio]{figures/exp_response_tomato_success.pdf}
\caption{TomatoHarvesting. Success rate (\%).}
\label{fig:response_response_tomato_success}
\end{subfigure}
\par\medskip
\begin{subfigure}[t]{0.495\textwidth}
\centering
\includegraphics[width=\linewidth,height=.20\textheight,keepaspectratio]{figures/exp_response_waste_queries.pdf}
\caption{WasteSorting. Query count.}
\label{fig:response_response_waste_queries}
\end{subfigure}
\hfill
\begin{subfigure}[t]{0.495\textwidth}
\centering
\includegraphics[width=\linewidth,height=.20\textheight,keepaspectratio]{figures/exp_response_tomato_queries.pdf}
\caption{TomatoHarvesting. Query count.}
\label{fig:response_response_tomato_queries}
\end{subfigure}
\caption{Comparison of oracle simulation and physical execution with VLM and human responses. Each response condition retains its own trial set.}
\label{fig:response_comparison}
\end{figure*}





\subsubsection{Ablation Study}~\label{exp:when_what_random}\label{exp:policy_ablation}

We conduct an ablation study to assess the contributions of query timing and query selection to task success and query efficiency. We compare seven configurations with oracle feedback in simulation. For concise reference, we use \textbf{W} to label query timing variants and \textbf{Q} to label query selection variants. With entropy gain selection fixed, we compare random (W1), conformal prediction (W2), and reward-based (W3) timing. With confidence timing fixed, we compare random (Q1) and reward-based (Q2) selection. Fully-Random randomizes both query timing and query selection, while Ours combines confidence timing with entropy gain selection. Table~\ref{tab:query_strategy_matrix} summarizes these configurations. W2 also changes the criterion for continuing a query episode, and W3 and Q2 use domain-specific discount factors for query evaluation and physical planning. Appendix~\ref{appx:ablation_configs} specifies these differences. Each configuration has $200$ trials per domain. Table~\ref{tab:policy_ablation} and Fig.~\ref{fig:policy_ablation_results} present the results. Appendix~\ref{appx:ablation_configs} provides the query policy configurations.

\begin{figure*}[!tp]
\centering
\begin{subfigure}[t]{0.495\textwidth}
\centering
\includegraphics[width=\linewidth,height=.20\textheight,keepaspectratio]{figures/exp_ablation_waste_success.pdf}
\caption{WasteSorting. Success rate (\%).}
\label{fig:ablation_ablation_waste_success}
\end{subfigure}
\hfill
\begin{subfigure}[t]{0.495\textwidth}
\centering
\includegraphics[width=\linewidth,height=.20\textheight,keepaspectratio]{figures/exp_ablation_tomato_success.pdf}
\caption{TomatoHarvesting. Success rate (\%).}
\label{fig:ablation_ablation_tomato_success}
\end{subfigure}
\par\medskip
\begin{subfigure}[t]{0.495\textwidth}
\centering
\includegraphics[width=\linewidth,height=.20\textheight,keepaspectratio]{figures/exp_ablation_waste_queries.pdf}
\caption{WasteSorting. Query count.}
\label{fig:ablation_ablation_waste_queries}
\end{subfigure}
\hfill
\begin{subfigure}[t]{0.495\textwidth}
\centering
\includegraphics[width=\linewidth,height=.20\textheight,keepaspectratio]{figures/exp_ablation_tomato_queries.pdf}
\caption{TomatoHarvesting. Query count.}
\label{fig:ablation_ablation_tomato_queries}
\end{subfigure}
\caption{Query timing and selection ablation. W1--W3 vary query timing and Q1--Q2 vary query selection. Confidence timing supports high success, and entropy gain selection reduces query count at similar success rates.}
\label{fig:policy_ablation_results}
\end{figure*}


\begin{table*}[t]
\centering
\caption{Query policy ablation. W1--W3 compare timing criteria and Q1--Q2 compare selection criteria.}
\label{tab:policy_ablation}
\label{tab:when_what_random}
\small
\setlength{\tabcolsep}{4pt}
\begin{tabular}{lllcccc}
\toprule
 & & & \multicolumn{2}{c}{WasteSorting} & \multicolumn{2}{c}{TomatoHarvesting} \\
\cmidrule(lr){4-5}\cmidrule(lr){6-7}
Condition & Timing & Selection & Success (\%) & Queries & Success (\%) & Queries \\
\midrule
Fully-Random & Random & Random & 73.50 & $10.25 \pm 6.11$ & 73.50 & $9.04 \pm 3.94$ \\
W1 & Random & Entropy Gain & 73.50 & $6.65 \pm 3.29$ & 73.50 & $7.43 \pm 2.72$ \\
W2 & Conformal Prediction & Entropy Gain & 74.00 & $7.09 \pm 2.58$ & 63.00 & $4.05 \pm 2.04$ \\
W3 & Reward & Entropy Gain & 99.00 & $17.65 \pm 2.89$ & 68.50 & $12.16 \pm 3.00$ \\
Q1 & Confidence & Random & 100.00 & $15.21 \pm 3.53$ & 97.50 & $12.68 \pm 2.53$ \\
Q2 & Confidence & Reward & 100.00 & $22.79 \pm 6.26$ & 97.50 & $15.25 \pm 2.95$ \\
Ours & Confidence & Entropy Gain & 100.00 & $7.83 \pm 2.16$ & 98.00 & $9.16 \pm 1.90$ \\
\bottomrule
\end{tabular}
\end{table*}

The timing comparison evaluates Fully-Random, W1--W3, and Ours. In WasteSorting, Fully-Random and W1 achieve $73.50\%$ success, and W2 achieves $74.00\%$. W3 achieves a similar success rate of $99.00\%$ compared with $100.00\%$ for Ours, but requires $17.65$ queries compared with $7.83$. Fully-Random, W1, and W2 use $10.25$, $6.65$, and $7.09$ queries respectively. In TomatoHarvesting, Fully-Random and W1 achieve $73.50\%$ success, while W2 and W3 achieve $63.00\%$ and $68.50\%$ compared with $98.00\%$ for Ours. Mean query counts are $9.04$ for Fully-Random, $7.43$ for W1, $4.05$ for W2, $12.16$ for W3, and $9.16$ for Ours.

These results show that task success depends on the query trigger as well as query selection. W1 replaces random selection in Fully-Random with entropy gain and reduces query count without improving success. Random timing can miss necessary verification because the trigger does not depend on current state uncertainty. W2 can also end queries after the action prediction set narrows to one candidate, even if the state facts required to execute the candidate action remain uncertain. W3 uses expected task reward to determine whether to initiate queries, then selects queries by entropy gain and repeats verification until confidence reaches $\tau$. In WasteSorting, object category verification directly informs the correct bin choice and task reward, so repeated verification may support the high success rate of $99.00\%$. However W3 requires $125.43\%$ more queries than Ours. In TomatoHarvesting, W3 uses $32.71\%$ more queries than Ours, but the success rate of W3 is $29.50$ percentage points lower. Therefore the reward-based configuration does not consistently provide the state verification needed for successful execution, even with more queries overall. Ours uses current state confidence as the query trigger and achieves high success in both domains with fewer queries than the reward-based configuration.

With confidence-based query timing fixed, different query selection strategies maintain high success rates but require different numbers of queries. The query selection comparison evaluates Fully-Random, Q1, Q2, and Ours. Fully-Random achieves $73.50\%$ success in both domains with $10.25$ queries in WasteSorting and $9.04$ in TomatoHarvesting. Q1 retains random selection but uses confidence timing, and success increases to $100.00\%$ and $97.50\%$ respectively. With confidence timing fixed, Q1, Q2, and Ours all achieve $100.00\%$ success in WasteSorting. Mean query counts are $15.21$ for Q1, $22.79$ for Q2, and $7.83$ for Ours. In TomatoHarvesting, Q1 and Q2 achieve $97.50\%$ success compared with $98.00\%$ for Ours. Mean query counts are $12.68$ for Q1, $15.25$ for Q2, and $9.16$ for Ours.

Confidence timing maintains $100.00\%$ success in WasteSorting and at least $97.50\%$ in TomatoHarvesting across all three query selection strategies. Under confidence timing, query selection primarily affects the number of queries needed to reach the required confidence. Random selection does not prioritize facts by their expected reduction in state uncertainty. Reward-based selection evaluates expected task reward without directly measuring the information provided by a query response. Entropy gain selects facts according to the expected reduction in belief entropy. A response can therefore eliminate more uncertainty per query, although normalized confidence need not increase after every response. Relative to Q1 and Q2, Ours reduces query count by $48.54\%$ and $65.64\%$ in WasteSorting and by $27.72\%$ and $39.92\%$ in TomatoHarvesting. Thus our query selection method reduces the number of queries required under confidence timing without reducing task success.

These results support the main contribution of our method through the distinct roles of query timing and query selection. Query timing determines whether the robot requests additional information before it executes actions under state uncertainty. Confidence timing achieves high success with random, reward-based, and entropy gain selection, while entropy gain under random timing reduces query count without improving success over Fully-Random. Once confidence timing establishes a high success level, entropy gain reduces the number of queries needed to reach the required confidence. Therefore our method combines a timing criterion that supports task success with a selection criterion that reduces query count at that success level.

\subsubsection{Threshold Experiments}~\label{exp:sys_threshold}

After the ablation comparison, we examine how the confidence threshold controls task success and query count. We vary $\tau$ from $0.0$ to $1.0$ in increments of $0.1$ with entropy gain selection fixed and correct oracle responses in simulation. Each threshold has $200$ trials per domain. Table~\ref{tab:threshold_results} and Fig.~\ref{fig:threshold_results} present the results. The results at $\tau=0.8$ use the same Ours trials as the baseline and ablation comparisons, as specified in Appendix~\ref{appx:oracle_reference}. The sweep characterizes the success and query-count trade-off without establishing a unique optimal threshold.

\begin{figure*}[!tp]
\centering
\begin{subfigure}[t]{0.49\textwidth}
\centering
\includegraphics[width=\linewidth]{figures/exp_threshold_success.pdf}
\caption{Task success.}
\label{fig:threshold_success}
\end{subfigure}
\hfill
\begin{subfigure}[t]{0.49\textwidth}
\centering
\includegraphics[width=\linewidth]{figures/exp_threshold_queries.pdf}
\caption{Mean query count.}
\label{fig:threshold_queries}
\end{subfigure}
\caption{Effect of the confidence threshold on task success and query count. The dotted line marks the default threshold of $\tau=0.8$. Requiring confidence of $1.0$ substantially increases query count with little or no additional success.}
\label{fig:threshold_results}
\end{figure*}


\begin{table*}[t]
\centering
\caption{Task success and query count across confidence thresholds.}
\label{tab:threshold_results}
\small
\begin{tabular}{ccccc}
\toprule
 & \multicolumn{2}{c}{WasteSorting} & \multicolumn{2}{c}{TomatoHarvesting} \\
\cmidrule(lr){2-3}\cmidrule(lr){4-5}
$\tau$ & Success (\%) & Queries & Success (\%) & Queries \\
\midrule
0.0 & 55.00 & $0.00 \pm 0.00$ & 52.00 & $0.00 \pm 0.00$ \\
0.1 & 54.00 & $0.69 \pm 0.84$ & 52.00 & $1.21 \pm 1.43$ \\
0.2 & 54.00 & $0.69 \pm 0.84$ & 52.50 & $1.27 \pm 1.42$ \\
0.3 & 54.00 & $0.69 \pm 0.84$ & 52.50 & $1.89 \pm 1.83$ \\
0.4 & 54.00 & $0.69 \pm 0.84$ & 88.50 & $5.42 \pm 2.30$ \\
0.5 & 54.00 & $0.69 \pm 0.84$ & 92.50 & $6.27 \pm 2.07$ \\
0.6 & 53.50 & $1.08 \pm 0.97$ & 92.50 & $7.20 \pm 1.99$ \\
0.7 & 100.00 & $7.74 \pm 2.25$ & 92.50 & $7.30 \pm 2.04$ \\
0.8 & 100.00 & $7.83 \pm 2.16$ & 98.00 & $9.16 \pm 1.90$ \\
0.9 & 100.00 & $7.74 \pm 2.25$ & 99.50 & $11.02 \pm 1.52$ \\
1.0 & 100.00 & $14.97 \pm 2.02$ & 100.00 & $19.00 \pm 1.43$ \\
\bottomrule
\end{tabular}
\end{table*}

In WasteSorting, success improves sharply at $\tau=0.7$ and remains constant at higher thresholds. With $\tau=0.0$, the robot makes no queries and achieves $55.00\%$ success. Thresholds from $0.1$ to $0.6$ yield success rates between $53.50\%$ and $54.00\%$ with at most $1.08$ queries on average. At $\tau=0.7$, success reaches $100.00\%$ with $7.74$ queries. Success remains $100.00\%$ over $\tau=0.7$--$0.9$, and query count stays between $7.74$ and $7.83$. However $\tau=1.0$ requires $14.97$ queries without further success improvement. The additional queries confirm that the strictest confidence requirement can increase interaction frequency after task success reaches its maximum.

In TomatoHarvesting, higher thresholds improve success over a wider range. Success is $52.00\%$ without queries at $\tau=0.0$ and remains at most $52.50\%$ through $\tau=0.3$. At $\tau=0.4$, success increases to $88.50\%$ with $5.42$ queries. Success reaches $92.50\%$ over $\tau=0.5$--$0.7$ and $98.00\%$ at $\tau=0.8$ with $9.16$ queries. Raising $\tau$ to $0.9$ increases success to $99.50\%$ with $11.02$ queries. At $\tau=1.0$, success reaches $100.00\%$ with $19.00$ queries. Therefore the increase from $0.9$ to $1.0$ adds $7.98$ queries for only $0.50$ percentage points of additional success.

The two domains require different confidence levels to reach high success, but both show limited gains from the strictest threshold. Relative to $\tau=0.8$, $\tau=1.0$ increases query count by $91.19\%$ in WasteSorting and $107.30\%$ in TomatoHarvesting. Success improves by $0.00$ and $2.00$ percentage points respectively. The common operating point $\tau=0.8$ balances success and query count because it achieves $100.00\%$ and $98.00\%$ success in the two domains with substantially fewer queries than $\tau=1.0$. Thus the threshold controls the amount of state verification, and the highest confidence requirement does not provide a proportional improvement in task success.

\subsubsection*{Interviews}\label{exp:interviews}

% Interview protocol and analyzed findings remain to be adds.
% Address comprehension of state verification and action selection questions.
% Report participant evidence about reasoning effort and repeated requests.

\subsubsection*{Discussion}\label{exp:discussion}

The experiments show that query timing and query selection make distinct contributions to successful task execution. Confidence timing maintains success rates close to $100\%$ across random, reward-based, and entropy gain selection in the oracle experiments. Entropy gain reduces the number of queries required at that success level. The Fully-Random and W1 comparison further shows that effective selection alone does not recover high success if query timing remains random. The threshold experiments also show that stricter state verification increases query count without a proportional improvement in success. Together, these results support the separation of a timing criterion that identifies the need for state information from a selection criterion that reduces the cost of obtaining sufficient information.

Physical execution with VLM feedback preserves the success advantage of Ours, but additional queries have different effects on execution time across domains. Ours uses more queries than KnowNo in both domains, while total time is longer in WasteSorting and shorter in TomatoHarvesting. In TomatoHarvesting, KnowNo uses nearly as many queries as Ours but executes longer plans and requires more interaction time. The failure analysis also identifies action choices that are feasible but do not advance the task. Action selection requires the VLM to reason about robot location and object states to determine the next action. Ours asks the VLM to verify current state facts and leaves action selection to the planner. These results compare complete systems with different action-selection mechanisms. The required judgments and subsequent action choices may help explain the execution costs, but the experiments do not measure reasoning effort directly.

% Provisional interview interpretation. This paragraph states hypotheses and must be revised using participant interview evidence.
The query format may also affect the effort required to provide human feedback. State verification gives participants a specific fact to assess from the current observations. Action selection requires participants to compare alternatives and consider how the next action contributes to the task goal. Therefore participants may find frequent state verification questions easier to answer than fewer action selection questions. Repeated verification can still interrupt the task or become burdensome if the purpose of each question is unclear. These possible responses suggest that question clarity and the reasoning required for each answer matter alongside query frequency. The participant interviews examine whether the differences in required judgment also appear in the perceived interaction effort.

% % Keep the experimental figures and tables before the conclusion.
% \clearpage

\input{section/06_conclusion}

\bmhead{Acknowledgements}
This research was supported in part by the National Research Foundation of Korea (NRF) grant funded by the Korea government (MSIT)(No. RS-2024-00461583) and the National Research Foundation of Korea (NRF) grant funded by the Korea government (MSIT) (No. RS-2024-00411007).

\bibliography{ref}

\begin{appendices}
\input{appendix/00_appx_content}
\section{Planning Algorithms}~\label{appx:algorithms}

This appendix provides the POMCP search procedures used in Sec.~\ref{method:PM}. The search budget is $n_{\mathrm{sim}}$, the maximum simulation depth is $d_{\max}$, the discount cutoff is $\epsilon$, and the UCB exploration constant is $c$. The search tree $T$ stores visit counts $N$, value estimates $V$, and sampled state particles $B(h)$ at history $h$.

\normalem
\subsection{POMCP Search with Applicable Actions}~\label{appx:alg_search}
Algorithm~\ref{alg:search} describes POMCP search with a fixed simulation budget. Each simulation samples a state from the particles stored at the root history. If the root contains no particles, the simulation starts from the committed knowledge state. These search-tree particles support action-value estimation and are separate from the weighted frontier used for confidence and query selection in Sec.~\ref{method:belief_rep}. The input $b$ is the implementation record containing the committed knowledge state and frontier data. SEARCH reads the committed state from this record and samples from $B(h)$, not from the weighted frontier. After verification, the executed action and committed successor select the retained observation branch. The branch supplies the root particles for the next search. Frontier weights do not replace those particles. After the simulations, the planner selects the highest-value action applicable to the committed state.

\begin{algorithm}[H]
\caption{POMCP Search with Applicable Actions}
\label{alg:search}
\textbf{Input}\quad belief $b$, simulation budget $n_{\mathrm{sim}}$ \\
\textbf{Output}\quad selected action $a \in \mathcal{A}$
\BlankLine

\Proc{\Search{$b$}}{
    $h \gets$ root history\\
    $s^{K} \gets knowledge(b)$\\
    
    \For{$i\gets1$ \KwTo $n_{\mathrm{sim}}$}{
        \eIf{$B(h) \neq \emptyset$}{
            sample $s \sim B(h)$ \label{alg_search:sample}\\
        }{
            $s \gets s^{K}$\\
        }
        
        \Simulate{$s, h, 0$} \label{alg_search:simulate}\\
    }
    $a \gets$ \SearchBest{$h, s^{K}, \textnormal{false}$}  \label{alg_search:search_best}\\
    \KwRet $a$ \\
}
\end{algorithm}

\subsection{POMCP Simulation with Applicable Actions}~\label{appx:alg_simulate}
Algorithm~\ref{alg:simulate} describes how simulations are constrained by symbolic action applicability. If a history node is expanded for the first time, only actions whose preconditions hold in the sampled symbolic state are inserted into the search tree. The generative model then samples the next state, observation, and reward while preserving the standard POMCP value-update structure.

\begin{algorithm}[H]
\caption{POMCP Simulation with Applicable Actions}
\label{alg:simulate}
\textbf{Input}\quad state $s$, history $h$, depth $depth$, maximum depth $d_{\max}$ \\
\textbf{Output}\quad simulated return $R$
\BlankLine

\Proc{\Simulate{$s, h, depth$}}{
    \If{$depth>0$ and $(depth\geq d_{\max}$ or $\gamma=0$ or $\gamma^{depth}<\epsilon)$}{
        \KwRet $0$\\
    }

    \If{there is no $a\llbracket s \rrbracket \in \mathcal{A}$}{
        $N(h) \gets N(h) + 1$\\
        \KwRet $0$\\
    }

    \If{history $h$ has no action children}{
        \ForEach{$a\llbracket s \rrbracket \in \mathcal{A}$ \label{alg_simulate:applicable}}{
            $T, N(ha), V(ha) \gets T \cup \{ha\}, 0, 0$ \label{alg_simulate:tree_expand}\\
        }
        $N(h) \gets N(h) + 1$\\
        $V(h)\gets$\Rollout{$s, depth$} \label{alg_simulate:rollout}\\
        \KwRet $V(h)$\\
    }

    $a \gets$ \SearchBest{$h, s, \textnormal{true}$} \label{alg_simulate:search_best}\\
    sample $(s^\prime, o, r)$ from simulator $\mathcal{G}(s, a)$ \label{alg_simulate:simulate}\\

    $R \gets r + \gamma \cdot$ \Simulate{$s^\prime, hao, depth + 1$}\\

    $B(h)\gets B(h) \cup \{s\}$\\
    $N(h) \gets N(h)+1$\\
    $N(ha) \gets N(ha)+1$\\
    $V(ha) \gets V(ha) + \frac{R - V(ha)}{N(ha)}$\\

    \KwRet $R$\\
}
\end{algorithm}

\subsection{Action Selection and Rollout}~\label{appx:alg_selection}
Algorithm~\ref{alg:selection} specifies the shared action-selection routine used during tree search and rollout. During simulation, unvisited applicable actions receive priority and visited actions are compared using UCB. The final root decision uses the current value estimate. Rollout actions are sampled uniformly from applicable actions. This keeps both planning and rollout consistent with the symbolic transition model.

\begin{algorithm}[H]
\caption{Action Selection and Rollout}
\label{alg:selection}
\textbf{Input}\quad history $h$, state $s$
\BlankLine

\Proc{\SearchBest{$h, s, use\_ucb$}}{
    $\mathcal{C} \gets \{a \mid a\llbracket s \rrbracket \in \mathcal{A},\ ha \in T\}$\\
    
    \If{$\mathcal{C} = \emptyset$}{
        \KwRet $\emptyset$\\
    }

    \eIf{$use\_ucb$}{
        \If{an action $a\in\mathcal{C}$ has $N(ha)=0$}{
            \KwRet an unvisited action in $\mathcal{C}$
        }
        \KwRet $\arg\max_{a \in \mathcal{C}} \left[V(ha) + c \sqrt{\frac{\log(\max\{1,N(h)\})}{N(ha)}}\right]$\\
    }{
        \KwRet $\arg\max_{a \in \mathcal{C}} V(ha)$\\
    }
}
    
\BlankLine

\Proc{\Rollout{$s, depth$}}{
    \If{$depth>0$ and $(depth\geq d_{\max}$ or $\gamma=0$ or $\gamma^{depth}<\epsilon)$}{
        \KwRet $0$\\
    }

    \If{there is no $a\llbracket s \rrbracket \in \mathcal{A}$}{
        \KwRet $0$\\
    }

    $a \sim \mathrm{Uniform}(\{ a \mid  a\llbracket s \rrbracket \in \mathcal{A}\}  )$ \label{alg_rollout:uniform}\\
    sample $s^\prime$ from $T(\cdot\mid s,a)$ and evaluate reward $r$ \label{alg_rollout:simulator}\\

    \KwRet $r + \gamma \cdot$ \Rollout{$s^\prime, depth + 1$}\\
}
\end{algorithm}

\input{appendix/06_comparison_algorithms.tex}
\section{Detailed Experimental Settings}~\label{appx:settings}

This appendix specifies the planner settings, query policies, domain models, and reporting conventions for Sec.~\ref{exp:setup}. The oracle experiments use simulated execution with correct responses. The VLM and human evaluations use physical robot execution. We aggregate trials within each domain for all reported comparisons.

\subsection{Planner and Evaluation Parameters}~\label{appx:planner_params}
Table~\ref{tab:app_planner_params} lists the default POMCP parameters for Ours. The oracle trial limit is $50$ physical actions, and the physical trial limit is $20$. Queries do not consume physical steps for Ours. Query-Action uses a limit of $50$ total decisions that includes queries and physical actions. The reward-based configurations use the discount factors specified below. These trial limits are not equivalent because a query consumes a decision for Query-Action but no physical step for Ours. Appendix~\ref{appx:oracle_budget_audit} reports the budget-limit failures.

\begin{table}[!htbp]
\centering
\caption{Planner and evaluation parameters used by Ours.}
\label{tab:app_planner_params}
\small
\begin{tabular}{lc}
\toprule
\textbf{Parameter} & \textbf{Value} \\
\midrule
POMCP simulations per planning step & $100$ \\
Maximum simulation depth & $20$ \\
Discount factor $\gamma$ & $0.2$ \\
UCB exploration constant $c$ & $1.0$ \\
Oracle physical action limit & $50$ \\
Physical action limit & $20$ \\
Default confidence threshold $\tau$ & $0.8$ \\
Discount cutoff $\epsilon$ & $0.005$ \\
Oracle frontier capacity & $8{,}000$ \\
Oracle search-node particle capacity & $8{,}000$ \\
Oracle transition sampling budget & $250$ \\
Oracle initial configurations per domain & $5$ \\
Trials per initial configuration and condition & $40$ \\
Oracle trials per domain and condition & $200$ \\
\bottomrule
\end{tabular}
\end{table}

With $\gamma=0.2$ and $\epsilon=0.005$, the discount cutoff stops simulation and rollout at depth $4$ because $0.2^4=0.0016<0.005$. Thus the nominal maximum depth of $20$ is not reached under the default settings. Root depth is zero, so a simulation can evaluate four transitions before the cutoff.

\subsection{Baseline Configurations}~\label{appx:baseline_configs}
The oracle baseline comparison evaluates KnowNo, IntroPlan, and Query-Action against Ours. The physical experiments compare KnowNo and Ours. Appendix~\ref{appx:baseline_algorithms} provides separate procedures for the three baselines. KnowNo and IntroPlan use GPT-4o for action generation and candidate scoring. This planning model is separate from Gemma 4 31B IT, which provides external feedback in the VLM experiments.

The recorded action generation temperature is $0.3$, and the candidate score temperature is $5.0$. The oracle TomatoHarvesting KnowNo batch does not record the generation temperature. KnowNo uses calibration thresholds $\hat{q}=0.845165$ for WasteSorting and $\hat{q}=0.8404$ for TomatoHarvesting in the oracle comparison. The physical VLM logs use $0.845165$ and $0.901119$ respectively, except for one WasteSorting trial with $\hat{q}=0.8704$. IntroPlan retrieves three examples and uses calibration thresholds $0.980493$ and $0.985552$ for the two domains in the same order. Query-Action uses the recorded domain-specific configuration. The discount factor is $0.9$ for WasteSorting and $0.5$ for TomatoHarvesting. Query cost is zero, the failure penalty is $10$, assumed response accuracy is $1.0$, and the planner uses $100$ simulations per decision.

\subsection{Ablation Configurations}~\label{appx:ablation_configs}
Fully-Random uses random timing and random selection. W1, W2, and W3 use entropy gain selection with random, conformal prediction, and reward-based timing respectively. Q1 and Q2 use confidence timing with random and reward-based selection respectively. Random timing initiates a query episode with probability $0.4$ per physical step. Random selection samples uniformly from facts that distinguish the current state hypotheses. Confidence timing initiates verification if belief confidence falls below $\tau=0.8$. After a trigger, Fully-Random, W1, Q1, Q2, and Ours continue verification until confidence reaches $\tau$ or no differentiating fact remains.

W2 initiates verification if the conformal action prediction set contains multiple candidates. The criterion is reevaluated after each response. The calibration thresholds are $0.8704$ for WasteSorting and $0.8404$ for TomatoHarvesting, with score temperature $5.0$. Four unparseable outputs count as failures in the success rate and contribute no successful-trial cost measurements.

W3 evaluates the expected cumulative reward of query and physical actions. A query episode begins if the highest query value exceeds the highest physical action value. After the initial trigger, entropy gain selects questions and confidence determines whether verification continues. Q2 starts queries using confidence and selects the fact with the highest query value. Q2 does not compare that value with physical action values to cancel a query. Both reward-based variants use zero query cost, a failure penalty of $10$, and assumed response accuracy $1.0$. The discount factors are $0.9$ for WasteSorting and $0.5$ for TomatoHarvesting and apply to query evaluation and physical action planning. The query evaluator uses at least $100$ simulations and increases the budget to the number of root candidates plus one if necessary. The remaining configurations use $\gamma=0.2$. Therefore W3 and Q2 change the discount factor of physical planning as well as the query criterion. W2 changes the continuation criterion as well as the initial query trigger. The comparisons evaluate these complete policy configurations, while Fully-Random versus W1 and Fully-Random versus Q1 retain the same physical planner settings.

\subsection{Oracle Reference Trials}~\label{appx:oracle_reference}
We use the complete Ours batch from the policy ablation as a shared oracle reference for the baseline and ablation comparisons as well as the threshold analysis at $\tau=0.8$. This reference contains $200$ trials per domain and is reused across these comparisons. The other threshold points retain their original trial batches.

% Parameter values below are verified against the current 02 domain YAML files.
% Historical per-run YAML snapshots are not included in the raw bundles.
\subsection{Domain Model Probabilities}~\label{appx:model_params}
Table~\ref{tab:app_model_params} lists the parameters in the domain configuration files. Transition probabilities determine symbolic successor hypotheses. Observation sampling generates simulated sensor responses, and observation likelihoods weight the represented hypotheses. Detection sampling accuracy and likelihood accuracy are specified separately. Physical execution supplies real sensor observations to the belief update.

\begin{table}[!htbp]
\centering
\caption{Domain model parameters for transition and observation models.}
\label{tab:app_model_params}
\small
\begin{tabular}{lll}
\toprule
\textbf{Domain} & \textbf{Model component} & \textbf{Probability} \\
\midrule
WasteSorting & Detect observed transition & $0.90$ \\
WasteSorting & Detect category transition & $0.50$ \\
WasteSorting & Pick transition & $0.90$ \\
WasteSorting & Place transition & $0.90$ \\
WasteSorting & Detect observation likelihood & $0.97$ \\
WasteSorting & Detect category observation & $0.85$ \\
WasteSorting & Pick/place observation & $0.95$ \\
WasteSorting & Detect sampling accuracy & $[0.94,1.00]$ \\
WasteSorting & False-positive likelihood & $0.03$ \\
\midrule
TomatoHarvesting & Navigate transition & $0.99$ \\
TomatoHarvesting & Detect transition & $0.95$ \\
TomatoHarvesting & Pick transition & $0.95$ \\
TomatoHarvesting & Scan transition & $0.90$ \\
TomatoHarvesting & Place/discard transition & $0.99$ \\
TomatoHarvesting & Detect observation likelihood & $0.90$ \\
TomatoHarvesting & Detect classification observation & $0.95$ \\
TomatoHarvesting & Scan observation & $0.85$ \\
TomatoHarvesting & Detect sampling accuracy & $[0.85,0.95]$ \\
TomatoHarvesting & False-positive probability & $0.03$ \\
TomatoHarvesting & Navigate observation & $0.95$ \\
TomatoHarvesting & Place observation & $1.00$ \\
\bottomrule
\end{tabular}
\end{table}

The WasteSorting category transition parameter $0.50$ creates competing category hypotheses after detection. Detection observation sampling uses the hidden object category, with uniform allocation of classification error across the other categories. Occluded objects and objects already held or placed in bins are excluded from simulated detection. In TomatoHarvesting, detection produces ripe or unripe labels, and rotten tomatoes share the ripe visual label. A subsequent scan distinguishes fresh from rotten tomatoes. The detect models combine individual object outcomes into a joint observation distribution. Detector confidence values do not directly scale the symbolic observation likelihood.

\subsection{Randomization and Episode Seeds}~\label{appx:randomization}
Each oracle trial records a pseudorandom seed. The runner uses the seed for Python \texttt{random}, NumPy sampling, initial arrangements, and randomized query decisions. Conditions within each oracle experiment share the seed assigned to the same domain, initial configuration, and repetition. Different policies can consume random numbers differently after different action choices. The shared Ours reference described in Appendix~\ref{appx:oracle_reference} retains its original seeds.

\subsection{Uncertainty Estimates and Time Measures}~\label{appx:metric_details}
Success rates use all attempted trials. In the main comparisons, query counts, query probabilities, plan lengths, and time measures use successful trials and are reported as mean $\pm$ sample standard deviation. Query probability is the fraction of physical steps that initiate at least one query. Several questions within one query episode increase query count but contribute only one triggered step. Success bars show $95\%$ Wilson confidence intervals. Cost error bars show sample standard deviations. A dash or the asterisk identified in the figure caption denotes a single successful trial for which sample standard deviation is unavailable.

Search time is the sum of directly recorded planning durations within a trial. The oracle plan-time comparison uses cumulative POMCP search time for Ours and Query-Action and cumulative candidate generation and scoring time for KnowNo and IntroPlan. Interaction time measures external requests and response processing and is reported for the physical experiments. Total time is the accumulated trial duration. Physical result tables and figures use the residual planning-time estimate
\[
t_{\mathrm{plan}}=t_{\mathrm{total}}-t_{\mathrm{interaction}}-t_{\mathrm{execution}}.
\]
This estimate includes overhead outside the separately recorded interaction and physical execution durations. Direct search time and residual planning time are distinct measurements.

\input{appendix/05_vlm_prompts.tex}

\section{Supplementary Evaluation Details}~\label{appx:exp}

This appendix specifies the threshold protocol and the failure categories used in Sec.~\ref{exp:sys_baseline}. The reported outcomes use the same trial sets as the corresponding main-text comparisons.

\subsection{Threshold Protocol}~\label{appx:threshold_exp}
The threshold study evaluates $\tau\in\{0.0,0.1,\ldots,1.0\}$ with entropy gain query selection and oracle responses. Each threshold has $200$ trials per domain. The planner parameters are those in Appendix~\ref{appx:planner_params}, and the shared reference at $\tau=0.8$ is specified in Appendix~\ref{appx:oracle_reference}. A threshold of zero suppresses queries because belief confidence cannot fall below zero. A threshold of one requires confidence to reach one before the robot continues, unless no differentiating fact remains. Table~\ref{tab:threshold_results} reports the success rates and successful-trial query counts.

\subsection{Ablation Reporting}~\label{appx:policy_comparison}
Table~\ref{tab:policy_ablation} contains Fully-Random, W1--W3, Q1--Q2, and Ours. Each configuration has $200$ attempted trials per domain. The success denominator includes the four unparseable W2 outputs. Cost metrics include successful trials only. The policy definitions and domain-specific discount factors are given in Appendix~\ref{appx:ablation_configs}.

\subsection{Oracle Trial Budget Audit}~\label{appx:oracle_budget_audit}
Query-Action counts both queries and physical actions toward the $50$-decision limit. Ours permits $50$ physical actions and does not subtract queries from that limit. Table~\ref{tab:app_query_action_budget} separates recorded plan failures from decision-limit termination. Every decision-limit trial contains exactly $50$ queries and physical actions combined. WasteSorting reaches the limit after $12$--$15$ physical actions and $35$--$38$ queries. TomatoHarvesting reaches the limit after $25$--$34$ physical actions and $16$--$25$ queries. The remaining failed trials terminate with the recorded label \texttt{PLAN FAILURE}. This label does not identify a specific missed verification.

\begin{table}[!htbp]
\centering
\caption{Termination categories of failed Query-Action oracle trials. Each domain has $200$ attempted trials.}
\label{tab:app_query_action_budget}
\small
\begin{tabular}{lrrr}
\toprule
Domain & Failures & Decision limit & Plan failure \\
\midrule
WasteSorting & $26$ & $7$ & $19$ \\
TomatoHarvesting & $85$ & $13$ & $72$ \\
\bottomrule
\end{tabular}
\end{table}

These counts identify trials affected by the decision limit but do not establish the outcome under a matched physical-action budget. The reported baseline results compare the recorded systems and limits. The effect of an equal physical-action budget requires a separate evaluation.

\subsection{Query Counts over All Attempted Oracle Trials}~\label{appx:all_trial_queries}
Table~\ref{tab:app_all_trial_queries} supplements the successful-trial means in the main text with query counts over all $200$ attempts per domain and method. Early failures can reduce query count because execution stops before task completion. Therefore the all-trial costs also require interpretation alongside success rate. Ours uses the shared reference in Appendix~\ref{appx:oracle_reference}.

\begin{table}[!htbp]
\centering
\caption{Oracle query counts over all attempted trials, including failed trials. Values are mean $\pm$ sample standard deviation.}
\label{tab:app_all_trial_queries}
\small
\begin{tabular}{lcc}
\toprule
Method & WasteSorting & TomatoHarvesting \\
\midrule
KnowNo & $1.92 \pm 1.40$ & $1.80 \pm 1.35$ \\
IntroPlan & $5.65 \pm 2.54$ & $6.91 \pm 2.98$ \\
Query-Action & $24.00 \pm 6.72$ & $9.73 \pm 5.73$ \\
Ours & $7.83 \pm 2.16$ & $9.15 \pm 1.89$ \\
\bottomrule
\end{tabular}
\end{table}

\subsection{VLM Failure Classification}~\label{appx:vlm_failure_analysis}
We classify failed physical trials using the available action candidates, external responses, and subsequent execution. A VLM failure occurs if a response incorrectly verifies a state fact or selects an action that causes failure despite an available feasible alternative. Feasible actions that do not advance the task can also cause a VLM failure if the resulting sequence exhausts the trial limit. A planner failure occurs if the candidate set contains only infeasible actions or the baseline autonomously selects an action whose preconditions are not satisfied. The termination label alone does not determine the category.

Table~\ref{tab:app_vlm_failures} reports the failure counts. In WasteSorting, Ours has three VLM feedback failures. KnowNo has three VLM failures and $17$ planner failures. In TomatoHarvesting, Ours has three VLM feedback failures. KnowNo has $18$ VLM failures and $12$ planner failures. Six of the $18$ VLM failures involve detection choices that exhaust the step budget despite feasible alternatives for task progress. The remaining $12$ involve incorrect action selections.

\begin{table}[!htbp]
\centering
\caption{Failure counts in the VLM evaluation. Percentages use failed trials as the denominator.}
\label{tab:app_vlm_failures}
\small
\begin{tabular}{llrrr}
\toprule
Domain & Method & Failed trials & VLM failures & Planner failures \\
\midrule
WasteSorting & KnowNo & $20$ & $3$ ($15\%$) & $17$ ($85\%$) \\
WasteSorting & Ours & $3$ & $3$ ($100\%$) & $0$ \\
TomatoHarvesting & KnowNo & $30$ & $18$ ($60\%$) & $12$ ($40\%$) \\
TomatoHarvesting & Ours & $3$ & $3$ ($100\%$) & $0$ \\
\bottomrule
\end{tabular}
\end{table}

\subsection{KnowNo Action Requests}~\label{appx:knowno_impl}
KnowNo generates candidate actions from domain information and the current symbolic state. Candidate scores and the calibrated threshold define the conformal prediction set. If the set contains multiple actions, the interface requests one action choice from the external expert. KnowNo requests one response per trigger. The selected action is then executed, and the execution result updates the state information used for the next decision. Appendix~\ref{appx:baseline_configs} lists the model and calibration parameters.

\input{appendix/04_user_study.tex}
\end{appendices}

\end{document}
